% II. RELATED WORK -- OWNERS: Conor and Emmett
% Cut for the page limit on 2026-10-23 (Patrick, with Claude Code): same sources and
% claims as Conor's and Emmett's version (commit 05287b0), with repetition removed.
\section{Related Work}\label{sec:related}
Faithfulness concerns whether an explanation reflects the model's actual reasoning, as
distinct from whether it appears plausible~\cite{jacovi2020}. Intervention studies probe
it by truncating, corrupting or paraphrasing chain-of-thought~\cite{lanham2023}, or by
adding biasing features that change answers without being mentioned in the
explanation~\cite{turpin2023}. Chen et al.\ extend hint tests to reasoning models,
measuring acknowledgement among responses that switch to the hinted
answer~\cite{chen2025}. Arcuschin et al.\ find unfaithful reasoning on naturally worded
prompts, including contradictory justifications and illogical
shortcuts~\cite{arcuschin2025}, so natural-setting evaluation is not novel here.
Parcalabescu and Frank distinguish self-consistency from faithfulness to internal
computation, comparing input contributions to answer and explanation with
CC-SHAP~\cite{parcalabescu2024}; our measure is narrower, checking agreement between
recoverable working and the final answer.

Formal approaches either generate the reasoning or verify it. Faithful CoT executes
generated symbolic instructions with a deterministic solver, making that chain
responsible for the answer without guaranteeing correct translation~\cite{lyu2023};
NSF-CoT combines formal verification with an LLM entailment component to assess written
steps~\cite{pramanik2026}. Our checker evaluates arithmetic the model already wrote and
does not produce its answer. Circuit tracing builds attribution graphs through
transcoder replacement models, using interventions to test hypotheses about
pathways~\cite{ameisen2025}. Deterministic verification therefore complements
interventions but cannot establish that written reasoning caused an answer.
